\section*{Your turn}

Check out tag \code{ch16} of the companion repository and run \code{npm ci}
once. Exercises 1 to 4 need nothing but the .NET SDK: they run a console
project and never open a socket. Exercises 5 to 8 need the database, the broker
and the seven services, and the compose line is
chapter~\ref{ch:identity-and-authorization}'s, unchanged. Exercises 6 and 7 each
rebuild the Catalog container, so leave them until you have finished the others
or expect Catalog to answer differently while they are in progress.

No service changes at this tag, so there is no new Postman collection either.
Exercise 8 uses \code{postman/mosaic-subgraphs}, which has been there since
chapter~\ref{ch:strangling-the-monolith}.

\begin{enumerate}
  \item \textbf{Run the five cases.} From the repository root:

  \begin{minted}[fontsize=\footnotesize,breaklines]{text}
dotnet run --project samples/executor-internals
  \end{minted}

  Five case names, each with its numbers and a \code{PASS}. Read the numbers
  before reading the code: four of the five are counts you can predict from this
  chapter, and one of them is a zero that took a section to explain.

  \item \textbf{Make a pure field cost something.} In the sample's
  \code{SampleTypes.cs}, change \code{Upper()} so that it returns a
  \code{Task<string>} and awaits \code{Task.Yield()} before returning, then run
  the first case alone:

  \begin{minted}[fontsize=\footnotesize,breaklines]{text}
dotnet run --project samples/executor-internals -- pure-fields-cost-no-task
  \end{minted}

  The case fails, and the number it reports tells you how many things changed.
  Work out why it is that number and not one larger. Put it back.

  \item \textbf{Find the sixty-fifth condition.} Write a small loop in the same
  project that builds a document selecting one field with
  \code{@skip(if: \$s1)}, then two with \code{\$s1} and \code{\$s2}, and so on,
  and execute each. Somewhere past sixty-four distinct variables the compiler
  stops being polite about it. Then change the loop so that all the fields share
  one variable and take it as far as you like, and confirm that the limit counts
  distinct conditions rather than annotated fields.

  \item \textbf{Watch the leader compile alone.} Raise the burst in the
  single-flight case from thirty-two to something that makes your machine work,
  and confirm the compilation count stays at one. Then make the document differ
  per request by interpolating the loop index into a field alias, and watch what
  the same burst costs when nothing can be coalesced.

  \item \textbf{Count the dashes.} Start the stack. Send the same document
  twice, because the first one is a cache miss and proves nothing, then read the
  last timeline line the service logged:

  \begin{minted}[fontsize=\footnotesize,breaklines]{text}
for i in 1 2; do
  curl -s -X POST http://localhost:5101/graphql \
      -H 'Content-Type: application/json' \
      -d '{"query":"{ __typename }"}' > /dev/null
done
docker logs mosaic-graph-mosaic-catalog-1 | grep 'parse ' | tail -1
  \end{minted}

  It should report a document cache hit; if it does not, you read the first
  request. Repeat for 5102 to 5107, changing the container name each time. Four
  services print \code{validate -} and three print a duration. Before you look it up, predict which three from what
  chapter~\ref{ch:identity-and-authorization} did, and check whether your list
  is the same as the one in that chapter's middleware count.

  \item \textbf{Make Catalog revalidate.} It has nothing to protect, which is
  the point. Add one line to the builder chain in its \code{Program.cs},
  together with the \code{using} that Accounts has and Catalog does not:

  \begin{minted}[fontsize=\footnotesize,breaklines]{csharp}
.AddMosaicAuthorization()
  \end{minted}

  Then rebuild that one container:

  \begin{minted}[fontsize=\footnotesize,breaklines]{text}
docker compose up -d --build mosaic-catalog
  \end{minted}

  Repeat exercise 5 for Catalog alone. The dash becomes a duration in a service
  with no \code{[Authorize]} anywhere in it. A build that fails on a missing
  extension method means the \code{using} did not go in. Put both lines back and
  rebuild before going on.

  \item \textbf{Give a Mosaic service a warmed cache.} Add a warmup task to
  Catalog that executes the storefront document with
  \code{MarkAsWarmupRequest()}, rebuild it, and send that document as your first
  request after startup. The timeline reports an operation cache hit on a
  request nothing had made before. Then decide whether you would ship it, and
  what you would have to know about your traffic first.

  \item \textbf{Verify it in Postman.} There is no collection for this chapter,
  so use the one that already talks to the subgraphs directly:

  \begin{minted}[fontsize=\footnotesize,breaklines]{text}
npx newman run postman/mosaic-subgraphs.postman_collection.json \
    -e postman/mosaic-subgraphs.local.postman_environment.json \
    --env-var jwtSecret="$(node -e "process.stdout.write(process.env.MOSAIC_JWT_SECRET||'')")"
  \end{minted}

  Run it twice without restarting anything, and read the Accounts and Catalog
  consoles either side of the second run. Every request in it is a document
  those services have already seen. Work out which of the two consoles gained
  validation time on the second pass and why that is not a property of the
  collection.
\end{enumerate}

Exercise 6 is the one to do properly. Most of this chapter is a fact about a
library, and would be true in any project that used it. That one is a fact about
yours: a line added for a reason that had nothing to do with performance, in a
service that gains nothing from it, reaching into a cache an earlier chapter
measured. There is no artefact anywhere in the build that could have told you,
because the schema does not change and neither does anything composed from it.
The only reason this book knows is that somebody counted requests on seven
consoles and noticed four of them looked different from the other three.
